% section: 漸化式を解く前に

\section{漸化式を解く前に}\label{節：漸化式を解く前に}
  \BookFigApproach

  本章で繰り返し使う問いは,「元の項より,次への変化が簡単になる量はないか」である.
  例えば$a_{n+1}=2a_n+3$では,定数$3$があるため等比型の式をそのまま使えない.
  両辺に$3$を加えると
  \[
    a_{n+1}+3=2(a_n+3)
  \]
  となり,\,$b_n=a_n+3$という同じ形の量が2倍されていると分かる.
  なぜ$3$を選ぶかは\sectionref*{節：特性方程式型}で導く.
  ここでは,変形の目標を具体的に持つことを確認しよう.
  \begin{itemize}
    \item $b_{n+1}-b_n=f(n)$なら,変化を足し戻す（階差型）.
    \item $b_{n+1}=f(n)b_n$なら,倍率を掛け合わせる（階比型）.
    \item $b_{n+1}=b_n$なら,その量は初項から変わらない.
    \item 既知の公式と同じ形なら,その公式の入力に当たる量を追う.
  \end{itemize}
  本書で「きれいな形」と呼ぶのは,こうした更新の仕組みを追いやすい形である.
  和や積の表示が得られた後,それをどこまで計算できるかは別に確認する.

  説明を読むときは,\textbf{観察と発想}と\textbf{式変形}を分けて追ってほしい.
  前者では,どの項を消すか,どの倍率をそろえるかを選ぶ.
  後者では,選んだ量を実際に代入し,初項から一般項を求め,元の量へ戻す.
  逆数・対数・割り算を使う前には,その操作が定義できるかも確かめる.
  候補を探す計算の途中で条件が必要だと分かった場合は,そこで条件を確認してから続ける.

  最初の数項を計算する実験も,変形の候補を見つける助けになる.
  数項の一致は予想の根拠であり,すべての項についての証明にはならない.
  予想した式が初期条件と元の更新式を満たせば,順に次の項が一意に決まる漸化式では,数学的帰納法で一致を示せる.
  \sectionref*{節：数学的帰納法}を参照してほしい.
  帰納法は,変形から得た式を確かめるときにも,実験から得た予想を証明するときにも使える.

  一般項を簡単な式で表すことに加えて,正値・極限・整数性などを調べる目標もある.
  \chapterref{章：実際の解き方}では方法の選択を練習し,発展演習と付録では目標の広がりを見る.
